\subsection{Multiplicities and finite extensions}

PROPOSITION 5. — Suppose A is semi-local and let \(\rho: A \to B\) be a ring homomorphism making B into a finitely generated A-module. Let N be a nonzero finitely generated B-module and let q be an ideal of A contained in the radical of A, and such that N/qN is of finite length. Among the maximal ideals of B (finite in number, by IV, § 2, no. 5, cor. 3 to prop. 9), denote by \(\mathfrak{m}_1, \ldots, \mathfrak{m}_r\) those for which we have \(\dim_{\mathrm{B}_{\mathfrak{m}_i}} (\mathrm{N}_{\mathfrak{m}_i}) = \dim_{\mathrm{B}} (\mathrm{N})\) . Put \(\mathbf{B}_i = \mathbf{B}_{\mathfrak{m}_i}\) and \(\mathfrak{q}_i = \mathfrak{q}\mathbf{B}_i\) for \(1 \leqslant i \leqslant r\) . Then we have the equalities

\[
\dim_ {A} (N) = \dim_ {B} (N),
\]

\[
e _ {\mathfrak {q} \mathrm{B}} ^ {\mathbf {B}} (\mathbf {N}) = \sum_ {i = 1} ^ {r} e _ {\mathfrak {q} _ {i}} ^ {\mathbf {B} _ {i}} (\mathbf {N} _ {\mathfrak {m} _ {i}}),
\]

\[
e _ {\mathfrak {q}} ^ {\mathrm{A}} (\mathrm{N}) = \sum_ {i = 1} ^ {r} \left[ \mathrm{B} / \mathfrak {m} _ {i}: \mathrm{A} / \rho^ {- 1} (\mathfrak {m} _ {i}) \right] e _ {\mathfrak {q} _ {i}} ^ {\mathrm{B} _ {i}} (\mathrm{N} _ {\mathfrak {m} _ {i}}).
\]

The first equality follows from § 2, no. 3, th. 1, c); the second follows from remark 4 of no. 1 (note that m \(_{i}\) belongs to V(qB) for \(1 \leqslant i \leqslant r\) since \(\rho^{-1}(m_{i}) \supset q\) by V, § 2, no. 1, prop. 1). Let us prove the third equality. Let E be a B-module of finite length; we have

\[
\operatorname{long} _ {\mathbf {A}} (\mathbf {E}) = \sum_ {m} \left[ \mathrm{B} / m: \mathrm{A} / \rho^ {- 1} (m) \right]. \operatorname{long} _ {\mathrm{B} _ {m}} \left(\mathrm{E} _ {m}\right),
\]

m ranging over the set of maximal ideals of B ; this is indeed evident when E is one of the B/m, and the general case follows from this, since E possesses a composition sequence whose quotients are isomorphic to B/m. Applying this formula to the B-modules N/q \(^{n+1}\) N, we deduce the desired equality by definition of multiplicities.

COROLLARY. — If \(\left[\mathrm{B}/\mathfrak{m}_{i}:\mathrm{A}/\rho^{-1}(\mathfrak{m}_{i})\right]=1\) for all i, we have \(e_{\mathfrak{q}}^{\mathrm{A}}(\mathrm{N})=e_{\mathfrak{q}\mathrm{B}}^{\mathrm{B}}(\mathrm{N})\) .

Lemma 1. — Let \(\rho: A \to B\) be a ring homomorphism, and let \(\mathfrak{p}\) be a prime ideal of A. Consider the following two properties:

\begin{enumerate}
\def\labelenumi{(\roman{enumi})}
\item
  the canonical homomorphism \(\tilde{\rho}\) from \(\mathbf{A}_{\mathfrak{p}}\) into \(\mathbf{A}_{\mathfrak{p}} \otimes_{\mathbf{A}} \mathbf{B}\) is bijective;
\item
  there exists a unique prime ideal \(\mathfrak{r}\) of B above \(\mathfrak{p}\) and the canonical homomorphism \(\rho_{\mathfrak{p}}\) from \(A_{\mathfrak{p}}\) into \(B_{\mathfrak{r}}\) is bijective.
\end{enumerate}

We have \((\mathrm{i}) \Rightarrow (\mathrm{ii})\). If \(\mathfrak{p}\) is minimal, or if \(\mathbf{B}\) is integral over \(\mathbf{A}\), we have \((\mathrm{i}) \Leftrightarrow (\mathrm{ii})^{1}\).

The ring \(\mathbf{A}_{\mathfrak{p}} \otimes_{\mathbf{A}} \mathbf{B}\) is identified with the ring of fractions \(S^{-1} \mathbf{B}\) of B defined by the multiplicative subset \(S = \rho(A - \mathfrak{p})\) of B. The prime ideals of \(S^{-1} \mathbf{B}\) are therefore the \(S^{-1} \mathfrak{q}\) , where \(\mathfrak{q}\) is a prime ideal of B such that \(\rho^{-1}(\mathfrak{q}) \subset \mathfrak{p}\) ; if \(\mathfrak{q}\) is such an ideal, \((S^{-1} \mathbf{B})_{S^{-1} \mathfrak{q}}\) is identified with \(\mathbf{B}_{\mathfrak{q}}\) (II, § 2, no. 5, prop. 11).

If condition (i) is satisfied, there exists (V, § 2, no. 1, lemma 1) a unique prime ideal \(\mathfrak{r}\) of B such that \(\rho^{-1}(\mathfrak{r}) = \mathfrak{p}\). Moreover, \(B_{\mathfrak{r}}\) is identified with the ring of fractions \((S^{-1}B)_{S^{-1}\mathfrak{r}}\), hence also with \((A_{\mathfrak{p}})_{s}\) where \(s\) is the inverse image of \(S^{-1}\mathfrak{r}\) under the isomorphism \(\tilde{\rho}: A_{\mathfrak{p}} \to S^{-1}B\); and \(\tilde{\rho}^{-1}(S^{-1}\mathfrak{r}) = (A - \mathfrak{p})^{-1}\mathfrak{p} = \mathfrak{p}A_{\mathfrak{p}}\), whence (ii).

Conversely, suppose (ii) is satisfied, and let r be the unique prime ideal of B lying over p. Since \((\mathrm{S}^{-1}\mathrm{B})_{\mathrm{S}^{-1}\mathrm{r}}\) is identified with \(B_{r}\) it suffices to prove that \(S^{-1}B\) is local with maximal ideal \(S^{-1}r\) , that is, that every prime ideal q of B such that \(\rho^{-1}(q) \subset p\) is contained in r. If p is minimal, we have \(\rho^{-1}(q) = p\) , hence q = r. If B is integral over A, there exists by V, § 2, no. 1, cor. 2 to th. 1, a prime ideal \(r'\) of B such that \(q \subset r'\) and \(\rho^{-1}(r') = p\) ; we necessarily have \(r' = r\) , hence \(q \subset r\) .

Lemma 2.--- Suppose A is semi-local; let q be an ideal of definition of A, and let \(\rho : A \to B\) be a ring homomorphism making B into a finitely generated A-module. Suppose that, for every prime ideal (necessarily minimal) \(\mathfrak{p}\) of A such that \(\dim(A/\mathfrak{p}) = \dim(A)\), there exists a unique prime ideal \(\mathfrak{r}\) of B lying over \(\mathfrak{p}\) and that the canonical homomorphism \(\rho_{\mathfrak{p}} : A_{\mathfrak{p}} \to B_{\mathfrak{r}}\) is bijective. Then we have \(\dim_{A}(B) = \dim(A)\) and \(e_{\mathfrak{q}}^{A}(B) = e_{\mathfrak{q}}^{A}(A)\). Let \(\mathfrak{S}_{A}\) (resp. \(\mathfrak{S}_{B}\)) be the set of prime ideals \(\mathfrak{p}\) of A such that

\[
\dim (\mathrm{A} / \mathfrak {p}) = \dim (\mathrm{A}) (\text { resp. } \dim_ {\mathrm{A}} (\mathrm{B} / \mathfrak {p} \mathrm{B}) = \dim_ {\mathrm{A}} (\mathrm{B}));
\]

we have \(\mathfrak{S}_{\mathbf{A}} \neq \emptyset\) . Let \(\mathfrak{p} \in \mathfrak{S}_{\mathbf{A}}\) ; there exists by hypothesis a prime ideal of B lying over \(\mathfrak{p}\) . Then one has \(\rho^{-1}(\mathfrak{pB}) = \mathfrak{p}\) (II, § 2, no. 5, cor. 3 to prop. 11), and

\[
\dim (\mathrm{A} / \mathfrak {p}) = \dim (\mathrm{B} / \mathfrak {p B}) = \dim_ {\mathrm{A}} (\mathrm{B} / \mathfrak {p B})
\]

by th. 1, b) and c) of § 2, no. 3. Consequently, we have

\[
\dim_ {\mathrm{A}} (\mathrm{B}) \geqslant \dim_ {\mathrm{A}} (\mathrm{B} / \mathfrak {p B}) = \dim (\mathrm{A} / \mathfrak {p}) = \dim (\mathrm{A}) \geqslant \dim_ {\mathrm{A}} (\mathrm{B}).
\]

This implies \(\mathfrak{S}_{\mathrm{A}} \subset \mathfrak{S}_{\mathrm{B}}\) and \(\dim(\mathrm{A}) = \dim_{\mathrm{A}}(\mathrm{B})\) . Conversely, if \(\mathfrak{p} \in \mathfrak{S}_{\mathrm{B}}\) , we have the inequalities

\[
\dim_ {A} (B / p B) = \dim_ {A} (B) = \dim (A) \geqslant \dim (A / p) \geqslant \dim_ {A} (B / p B),
\]

whence \(p \in \mathfrak{S}_A\) and \(\mathfrak{S}_B = \mathfrak{S}_A\) . By prop. 3 of no. 1 and its corollary, we have

\[
e _ {\mathfrak {q}} ^ {\mathrm{A}} (\mathrm{A}) = \sum_ {\mathfrak {p} \in \mathfrak {S} _ {\mathrm{A}}} e _ {\mathfrak {q}} ^ {\mathrm{A}} (\mathrm{A} / \mathfrak {p}) \quad \text { and } \quad e _ {\mathfrak {q}} ^ {\mathrm{A}} (\mathrm{B}) = \sum_ {\mathfrak {p} \in \mathfrak {S} _ {\mathrm{B}}} \operatorname{long} _ {\mathrm{A} _ {\mathfrak {p}}} (\mathrm{A} _ {\mathfrak {p}} \otimes_ {\mathrm{A}} \mathrm{B}) e _ {\mathfrak {q}} ^ {\mathrm{A}} (\mathrm{A} / \mathfrak {p});
\]

by Lemma 1, we have \(\operatorname{long}_{A_{\mathfrak{p}}}(A_{\mathfrak{p}} \otimes_{A} B) = 1\) for every \(\mathfrak{p} \in \mathfrak{S}_{A}\), whence \(e_{\mathfrak{q}}^{A}(A) = e_{\mathfrak{q}}^{A}(B)\) .

PROPOSITION 6. --- Suppose A is semi-local and reduced; let q be an ideal of definition of A; let A' be the total ring of fractions of A, and let B be a finite sub-A-algebra of A'. Then B is semi-local and qB is an ideal of definition of it. Suppose that, for every maximal ideal m of B such that \(\dim(\mathbf{B}_{\mathfrak{m}})=\dim(\mathbf{B})\) , one has \([\mathbf{B}/\mathfrak{m}:\mathbf{A}/(\mathbf{A}\cap\mathfrak{m})]=1\) . Then one has \(e_{\mathfrak{q}}^{\mathbf{A}}(\mathbf{A})=e_{\mathfrak{q}\mathbf{B}}^{\mathbf{B}}(\mathbf{B})\) .

By IV, § 2, no. 5, cor. 3 of prop. 9, B is semi-local with ideal of definition qB. We have \(e_{qB}^{B}(B) = e_{q}^{A}(B)\) by the corollary to prop. 5. Since A' is identified with \(\prod_{p} A_p\) where p runs over the set of minimal prime ideals of A (IV, § 2, no. 5, prop. 10), the canonical map \(A_{p} \to A_{p} \otimes_{A} B\) is bijective for every minimal prime ideal p of A. It then follows from lemmas 1 and 2 that \(e_{q}^{A}(B) = e_{q}^{A}(A)\) , whence the proposition.

Example. --- Let \(k\) be a field of characteristic \(\neq 2\) and take for A the local ring \(k[[X, Y]]/(X^{2} + Y^{2})\) with residue field \(k\) . Take \(B = k[[X, T]]/(T^{2} + 1)\) where \(T = Y/X\) . Distinguish two cases: if -1 is the square of an element \(i\) of \(k\) , B has two maximal ideals generated respectively by \(\{X, T + i\}\) and \(\{X, T - i\}\) , they have residue field \(k\) , and one has \(e_{\mathfrak{m}_{A}}^{A}(A) = e_{\mathfrak{m}_{A}B}^{B}(B) = 2\) . If -1 is not a square in \(k\) , B has a unique maximal ideal (X) with residue field \(k[T]/(T^{2} + 1)\) , and one has \(e_{\mathfrak{m}_{A}}^{A}(A) = 2\) , \(e_{\mathfrak{m}_{A}B}^{B}(B) = 1\) .

